%% ma: L10-B14-0080
%% lop: 10
%% chuong: 5
%% bai: 14
%% dang: 10.14.13
%% loai: TN4
%% muc: VD
%% dap_an: "C"
%% nguon: "chuyen-de-cac-so-dac-trung-cua-mau-so-lieu-khong-ghep-nhom-toan-10-chuong-trinh-moi.pdf (GV. Phan Nhật Linh) (P05), Bài 03 Dạng 2 Phần I Câu 27"
%% kien_thuc: "Lớp 10 Bài 13 (số trung bình), Bài 14 (phương sai, độ lệch chuẩn)"
%% trang_thai: chinh-thuc
%% ghi_chu: "Đồ họa bản đồ gốc (TTXVN) thay bằng biểu đồ cột TikZ cùng số liệu. Tự duyệt theo yêu cầu giáo viên 22/10/2026."
\begin{ex}
Lượng khách quốc tế đến Việt Nam trong $6$ tháng đầu năm 2022 theo châu lục (đơn vị: nghìn lượt khách, nguồn: Tổng cục Thống kê) được cho trong biểu đồ sau:
\begin{center}
\begin{tikzpicture}[x=1.6cm,y=.0085cm]
  \foreach \v/\t [count=\i] in {392.1/Châu Á,96.4/Châu Âu,81.0/Châu Mỹ,30.0/Châu Đại Dương,2.4/Châu Phi}{
    \fill[teal!45] (\i-.3,0) rectangle (\i+.3,\v);
    \draw (\i-.3,0) rectangle (\i+.3,\v);
    \node[above,font=\scriptsize] at (\i,\v) {\pgfmathprintnumber[use comma,fixed,precision=1,fixed zerofill]{\v}};
    \node[below,font=\scriptsize,align=center,text width=1.4cm] at (\i,0) {\t};
  }
  \draw[truc] (0.4,0)--(5.7,0);
  \draw[truc] (0.4,0)--(0.4,450);
  \foreach \y in {100,200,300,400}{\draw (0.35,\y)--(0.45,\y); \node[left,font=\scriptsize] at (0.35,\y) {$\y$};}
\end{tikzpicture}
\end{center}
Phương sai của mẫu số liệu nhận được là
\choice{$1609{,}016$}{$29\,609{,}0816$}{\True $19\,609{,}0816$}{$19\,600{,}0816$}
\loigiai{Mẫu số liệu: $392{,}1$; $96{,}4$; $81{,}0$; $2{,}4$; $30{,}0$. $\overline x=\dfrac{601{,}9}5=120{,}38$.

$s^2=\dfrac15\sum\left(x_i-120{,}38\right)^2=19\,609{,}0816$. Chọn C.}
\end{ex}
